\section{Integral annihilators and reflection for the four pure quadratic labels}
\label{sec:abelian-reflection-strict-vanishing}

\paragraph{Fixed field and original coefficients.}
Put $p=269$, $K=\mathbb Q(\omega)$, and
$H=\mathbb Q(\zeta_{24p})$.
The field $H$ is an auxiliary residual splitting field; it does
not replace any of the prescribed growing hosts in (FC.1).
For $d\in\{1,-4,8,-8\}$ put
$\delta_d=\chi_d|_{G_K}$, with $\chi_1=1$.
Let $c_p$ be the actual cyclotomic character and let $\omega_p$
be its Teichmuller part.  These conventions are used both in
the cohomology and in the character evaluations below.

For a discrete coefficient module $D$ use the full strict group
\[
 \begin{gathered}
 S_{\rm str}(K,D)=\ker\bigl(H^1_{\rm cts}(G_K,D)
                                  \longrightarrow\mathcal L(D)\bigr),\\
 \mathcal L(D)=\prod_{v\mid p}H^1(G_{K_v},D)
                   \times\prod_{v\nmid p,\infty}H^1(I_v,D).
 \end{gathered}
 \tag{G168.1}
\]
No condition at $p$ is weakened to an inertia condition.

\begin{theorem}[AC-G168: broad original-coefficient strict vanishing]
\label{AC-G168}
Define subsets of $\{1,\ldots,267\}$ by
\[
 \begin{aligned}
 \mathcal R_1=\mathcal R_{-4}=\mathcal R_8
      &=\{1,\ldots,267\},\\
 \mathcal R_{-8}
      &=\{1,\ldots,267\}\setminus\{89,92,99,168,175,178\}.
 \end{aligned}
 \tag{G168.2}
\]
Let $E/\mathbb Q_p$ be finite with its original ring of integers
$O$, allowing arbitrary absolute ramification.
Let $m\in\mathbb Z$, $e\geq1$, and let
$\kappa:G_K\to1+\mathfrak m_O$ be a continuous character
unramified at every finite place away from $p$.
If $r\equiv e+m\pmod{268}$ belongs to $\mathcal R_d$, then
\[
 S_{\rm str}\bigl(K,(E/O)
       (\delta_d\omega_p^{-m}\kappa^{-1}c_p^{-e})\bigr)=0.
 \tag{G168.3}
\]
The analogous residual strict groups vanish over every finite
field of characteristic $p$ for these residues and also for
$r=0$.  The latter residual statement is not used to claim
the original-coefficient assertion (G168.3) at $r=0$.
There are respectively $267,267,267,261$ residue classes in
(G168.2).  No assertion of nonvanishing is made at the six
uncovered residues for $\delta_{-8}$.

For $m=0$, $\kappa=1$, the four original finite labels are
exactly $1,\chi_{-4}|_{G_K},\chi_8|_{G_K},\chi_{-8}|_{G_K}$.
The theorem does not assert that their positive derived indices
are zero or supply the separate actual-window comparison.
\end{theorem}

\begin{proof}
\proofstep{1. Complete character dictionary of the auxiliary field.}
Write $G=\operatorname{Gal}(H/\mathbb Q)$ and
$J=\operatorname{Gal}(H/K)$.
Then
\[
 |G|=\varphi(24)\varphi(p)=2144,
       \qquad |J|=1072,
       \qquad p\nmid |G|.
 \tag{G168.4}
\]
The exponent of $G$ is $268=p-1$.
Consequently every character of $G$ takes its values in
$\mathbb Z_p^\times$, and its reduction is a character over
$\mathbb F_p$ with the same order.
Let $\omega$ denote the Teichmuller lift on $G$ of its action
on $\mu_p$.  On the cyclotomic automorphism $\sigma_a$ it
reduces to $a\pmod p$.
All characters of $G$ are uniquely
\[
 \chi_{D,k}=\chi_D\omega^k,\qquad
 D\in\{1,-4,8,-8,-3,12,-24,24\},\quad0\leq k<268.
 \tag{G168.5}
\]
These eight $\chi_D$ are the eight quadratic characters of
$(\mathbb Z/24\mathbb Z)^\times$.
Their parity is the sign of $D$, with $D=1$ even.
Thus exactly $1072$ of (G168.5) are odd.
For an odd character the primitive conductor is
\[
          f_{D,k}=|D|\quad(k=0),\qquad
          f_{D,k}=|D|p\quad(k\ne0).
 \tag{G168.6}
\]
In particular every such conductor divides $24p$.
The coprime conductor factors and the primitivity of a nontrivial
character modulo the prime $p$ justify (G168.6), including
all imprimitive characters of the ambient field.

For a finite $p$-primary $\mathbb Z_p[G]$-module $A$ write
$A(\chi)=e_\chi A$, where
\[
                  e_\chi=|G|^{-1}\sum_{g\in G}\chi(g)^{-1}g.
 \tag{G168.7}
\]
These idempotents are integral over $\mathbb Z_p$ and give an
exact direct-sum decomposition.
Put $A_H=\operatorname{Cl}(H)\otimes\mathbb Z_p$.

\proofstep{2. Integral odd-character annihilators at primitive conductors.}
For $F=\mathbb Q(\zeta_f)$ use the literal convention
\[
 \theta_f=\frac1f\sum_{\substack{1\leq a<f\\(a,f)=1}}
                a\sigma_a^{-1},\qquad
 I_f=\mathbb Z[\operatorname{Gal}(F/\mathbb Q)]\theta_f
                        \cap\mathbb Z[\operatorname{Gal}(F/\mathbb Q)].
 \tag{G168.8}
\]
Sharifi, Definition~4.5.1, Lemma~4.5.2 and Theorem~4.5.6,
give $I_f\operatorname{Cl}(F)=0$.
Their complete applied printed proofs were read, with the
upstream class-field and Gauss-sum inputs retained.
For every positive $c$ prime to $f$,
\[
 S_{f,c}=(c-\sigma_c)\theta_f
      =\sum_{\substack{1\leq a<f\\(a,f)=1}}
          \left\lfloor\frac{ca}{f}\right\rfloor\sigma_a^{-1}
                 \in I_f.
 \tag{G168.9}
\]
Indeed the coefficient of $\sigma_a^{-1}$ in
$\sigma_c\theta_f$ is $\{ca/f\}$; subtraction proves both
the equality and integrality without division in $\mathbb F_p$.
The source uses the negative of this generator, which belongs
to the same annihilating ideal.

For $\chi_{D,k}$ at its primitive conductor, the exact reduction
of the eigenvalue of (G168.9) is
\[
 s_{D,k}(c)=
 \sum_{\substack{1\leq a<f_{D,k}\\(a,f_{D,k})=1}}
       \left\lfloor\frac{ca}{f_{D,k}}\right\rfloor
                    \chi_D(a)a^{-k}\pmod p.
 \tag{G168.10}
\]
When $k=0$ no inverse modulo $p$ is needed; for $k\ne0$,
the units in the sum are prime to $p$.
The inverse exponent in (G168.10) comes from
$\sigma_a^{-1}$, not from an implicit change of the original
coefficient character.
A nonzero value makes the exact $p$-adic eigenvalue a unit,
and hence kills the corresponding $p$-class component of $F$.
This applies also to the character $\omega$: the integral
floor sum is used before reduction, so no invalid cancellation
of a Bernoulli denominator is involved.

The finite certificate enumerates all $1072$ odd pairs $(D,k)$
and every $2\leq c<32$ prime to the relevant $f_{D,k}$.
There are $16750$ such listed evaluations.
Exactly $1069$ rows have a displayed nonzero witness.
The only rows without one in this finite list are
\[
                   (-8,100),\qquad(24,93),\qquad(24,179).
 \tag{G168.11}
\]
The root implementation uses exact integer floors and modular
powers.  A separate peer implementation, using elementary
quadratic residue rules, grouping by residues modulo $269$,
and Horner evaluation along the primitive root $2$, reproduces
all $16750$ values.  The exact two implementation hashes and
input/output hashes are retained with this theorem.
The three zero rows are not a class-nonvanishing assertion and
are not a claim that every integral annihilator is zero there.

Here is the required descent from primitive conductor to $H$.
For $F=\mathbb Q(\zeta_{f_{D,k}})\subseteq H$, put
$U=\operatorname{Gal}(H/F)$.
Both norm/extension identities are exact:
\[
 N_{H/F}j_{H/F}=|U|,\qquad
             j_{H/F}N_{H/F}=\sum_{u\in U}u.
 \tag{G168.12}
\]
Since $|U|$ is prime to $p$, extension identifies
$A_F$ with $A_H^U$.
For surjectivity, if $x\in A_H^U$, then
$j_{H/F}N_{H/F}x=|U|x$, and multiplication by $|U|$ is
invertible on this finite $p$-group.
The identification is equivariant for the remaining Galois
action.  Since $\chi_{D,k}$ is trivial on $U$, its component
is therefore unchanged by this descent.
This also follows from the complete proof of Sharifi
Lemma~1.3.5, read at the source.
Thus no spurious Euler factor from a larger modulus is imposed.
We have proved
\[
 A_H(\chi_{D,k})=0\quad\hbox{for every odd pair outside (G168.11).}
 \tag{G168.13}
\]

\proofstep{3. The exact reflection direction, including its correction.}
For clarity we prove the part of Spiegelung actually used here.
Sharifi Theorem~4.4.2, whose entire statement and printed proof
were read, states for an odd character $\chi$ that
\[
 r_p A_H(\omega\chi^{-1})-\delta_\chi
                              \leq r_p A_H(\chi).
 \tag{G168.14}
\]
The correction can be nonzero only if $\chi=\omega$ and
$H(\mu(H)^{1/p})/H$ is unramified, in which case it is one.
The class groups in this statement are ordinary class groups.
No group with primes above $p$ inverted is substituted into it.

Here is its full specialized Kummer argument.
Let $B\subset H^\times$ consist of elements whose $p$th
roots generate unramified extensions of $H$.
Class-field theory and Kummer theory give the canonical
$G$-equivariant exact sequence
\[
 0\longrightarrow
       (B\cap\mathcal O_H^\times)/\mathcal O_H^{\times p}
 \longrightarrow\operatorname{Hom}(A_H/pA_H,\mu_p)
 \longrightarrow A_H[p].
 \tag{G168.15}
\]
There is no surjectivity assertion at the last term.
To see the map, an unramified Kummer radical $x$ has every
valuation divisible by $p$, so $(x)=\mathfrak a^p$; send it
to $[\mathfrak a]$.
The kernel consists exactly of unit radicals modulo $p$th
powers.  The unramified class-field pairing identifies the
middle term.  This is the sequence in Sharifi Lemma~1.4.6;
its entire proof and the applied Proposition~1.4.4 were read.

Let $j$ be complex conjugation in $G$.
The minus part of the unit group modulo torsion is zero:
for a CM field its units have the same rank as the units of
the maximal real subfield, and the anti-invariant units are
roots of unity.  This is Sharifi Lemma~1.4.12, with its complete
short proof read; equivalently it follows from the
equivariant Dirichlet unit theorem.
Since $2$ is invertible modulo $p$, the only possible odd
unit-character contribution in (G168.15) is the line
$\mu(H)/\mu(H)^p$, of character $\omega$.

The $p$-primary roots of unity in $H$ are exactly $\mu_p$:
$\mu_{p^2}$ cannot lie in an extension of $\mathbb Q$ of
degree prime to $p$.
At any place above $p$ the field $H$ is the compositum of
$\mathbb Q_p(\mu_p)$ with the unramified quadratic extension,
since $p\equiv5\pmod{24}$ has order two modulo $24$.
Adjoining $\mu_{p^2}$ is totally ramified of degree $p$ over
this local field.  For example the usual Eisenstein
cyclotomic polynomials after the substitution $X=1+T$ show
ramification degrees $p-1$ and $p(p-1)$; unramified base
change preserves their ratio.
Thus every nonzero class of $\mu(H)/\mu(H)^p$ gives a
ramified Kummer extension at $p$.
The odd kernel in (G168.15) is zero, including its
$\omega$ component.  In particular $\delta_\chi=0$ in
(G168.14).
Although $H$ has four places above $p$, there is no further
local correction in this ordinary-class argument.

For an even character $\nu$, put $\chi=\omega\nu^{-1}$.
The action on a homomorphism to $\mu_p$ is
\[
                    (g f)(x)=g\,f(g^{-1}x).
 \tag{G168.16}
\]
Consequently the dual of the $\nu$ component of $A_H/pA_H$
has character $\omega\nu^{-1}$, and (G168.15) injects that
dual into $A_H[p](\chi)$.
This proves precisely
\[
 r_p A_H(\nu)\leq r_p A_H(\omega\nu^{-1}),\qquad
 A_H(\omega\nu^{-1})=0\ \Longrightarrow\ A_H(\nu)=0.
 \tag{G168.17}
\]
For the zero conclusion, a finite $p$-group with zero quotient
modulo $p$ is zero.  No canonical identification of $A_H[p]$
with $A_H/pA_H$ is assumed.

Combining (G168.13) and (G168.17), all character components
of $A_H$ vanish except possibly the following six:
\[
 \begin{gathered}
 (-8,100),\ (24,93),\ (24,179),\\
 (-8,169),\ (24,176),\ (24,90).
 \end{gathered}
 \tag{G168.18}
\]
The second line consists of the reflected pairs
$(D,1-k\bmod268)$ from the first line.
This is an upper support statement, not a nonvanishing or
class-number computation for the six possible components.

\proofstep{4. Strict residual classes and the two rational extensions.}
Let $k$ be any finite field of characteristic $p$, and put
$a=\bar\delta_d\bar\omega_p^{-r}$ on $G_K$.
The field $H$ splits $a$, and $[H:K]=1072$ is prime to $p$.
Restriction therefore injects $H^1(G_K,k(a))$ into
$H^1(G_H,k)$ with its twisted quotient action.
For a cocycle $z$ the restriction is a homomorphism and
\[
             z(\tau g\tau^{-1})=a(\tau)z(g)
                       \quad(g\in G_H,\ \tau\in G_K).
 \tag{G168.19}
\]
This follows directly from the cocycle identity and $a|_{G_H}=1$.

If the original class is strict, this homomorphism is zero
on every inertia subgroup of $G_H$ and on every full
decomposition group above $p$.
Indeed the corresponding local class is a coboundary, and
on the subgroup where the coefficient is trivial that
coboundary is literally zero.
It therefore factors through the class group with the primes
above $p$ removed.  As that group is a quotient of the ordinary
class group, there is an injection
\[
 S_{\rm str}(K,k(a))\lhook\joinrel\longrightarrow
        \operatorname{Hom}_{\mathbb F_p[J]}
                        (A_H/pA_H,k(a)).
 \tag{G168.20}
\]
Only this injection is needed; no equality with the ordinary
class-group Hom space is asserted.
The Artin map is Galois equivariant for the ordinary conjugation
action, as in Sharifi Proposition~1.3.2.
Thus (G168.19) selects the class character $a$, not $a^{-1}$.
This is separate from the cyclotomic duality in (G168.16).

The two extensions to $G$ of $a|_J$ are
\[
           \chi_d\omega^{-r},\qquad
                         \chi_d\chi_{-3}\omega^{-r}.
 \tag{G168.21}
\]
One is odd and the other even.
Exact idempotents show that only these two class components
can contribute to (G168.20), also after scalar extension to $k$.
The rational quadratic partners are
\[
 \begin{array}{c|rrrr}
 d&1&-4&8&-8\\ \hline
 d\text{ after multiplication by }\chi_{-3}&-3&12&-24&24.
 \end{array}
 \tag{G168.22}
\]
These identities were independently checked on every unit
modulo $24$ in the residue-translation certificate.

\proofstep{5. Exact good residues, without a parity shortcut.}
For an odd component $(D,k)$ use its row in the integral
table; for an even component use the reflected odd row
$(D,1-k\bmod268)$.
For $d=1,-4,8$, both discriminants in (G168.22) are absent
from the possible exceptions (G168.18), so both selected
class components vanish for every $r\pmod{268}$.

For $d=-8$ the complete list of possible failures is
\[
 \begin{array}{c|c|c}
 r&\text{uncovered odd row}&\text{use}\\ \hline
 89 &(24,179)&\text{direct}\\
 92 &(24,93)&\text{reflected}\\
 99 &(-8,100)&\text{reflected}\\
 168&(-8,100)&\text{direct}\\
 175&(24,93)&\text{direct}\\
 178&(24,179)&\text{reflected}.
 \end{array}
 \tag{G168.23}
\]
For example, direct use has exponent $k=-r$, whereas the
reflected row has exponent $1+r$.
The table records every $0\leq r<268$, both selected
characters, their correct odd source rows, and a displayed
nonzero integral operator witness whenever available.
This is an independent translation of the root and peer
annihilator table, not a re-execution of its floor sums.
It proves the residual assertions, including $r=0$ for all
four labels.  The six uncovered residues are left open.

\proofstep{6. Lifting over the original, arbitrarily ramified ring.}
Let $\varpi$ be a uniformizer of the original $O$, and set
\[
       D=(E/O)(\delta_d\omega_p^{-m}\kappa^{-1}c_p^{-e}).
 \tag{G168.24}
\]
It is divisible and $D[\varpi]\simeq k(a)$, with
$r\equiv e+m\pmod{268}$.
If $r\ne0$, the inertia character at $p$ has nontrivial
reduction $\bar c_p^{-r}$, because $K_p/\mathbb Q_p$ is
unramified and $\delta_d$ is unramified at $p$.
For an appropriate inertia element the scalar minus one is
an $O$-unit.  Hence $D^{G_{K_p}}=0$.

Away from $p$, the cyclotomic, Teichmuller and $\kappa$
characters are unramified.  The remaining inertia action
has order at most two, so its fixed module is either zero
or the whole divisible module $D$.
For every local group $L$ used in (G168.1), therefore,
\[
                         D^L/\varpi D^L=0.
 \tag{G168.25}
\]
The long exact sequence for multiplication by $\varpi$
gives injectivity $H^1(L,D[\varpi])\to H^1(L,D)$.
A strict global class killed by $\varpi$ lifts to a class
in $H^1(G_K,D[\varpi])$; local injectivity makes the lift
strict.  It vanishes by the residual result already proved.
Every continuous cocycle from compact $G_K$ to discrete
$E/O$ has finite image and is killed by a power of $\varpi$.
A nonzero strict class would therefore yield a nonzero
one killed by $\varpi$, which is impossible.
This proves (G168.3) on the original coefficient module.

The exclusion of $r=0$ from this lifting argument is deliberate.
Each rational quadratic $\chi_d$ is unramified at $p$;
on $G_{K_p}$ its Frobenius value is the square of its
rational value, hence is one.
Thus all four local residual characters are trivial at
$r=0$, including $\chi_8$ and $\chi_{-8}$.
Local fixed modules for the full $p$-adic character can then
fail to be divisible, and the injection used in (G168.25)
does not follow.  No such lift is inferred from residual
strict vanishing alone.
\end{proof}

\paragraph{Proof and source boundaries.}
The full applied primary passages of Sharifi's Chapter~4
Theorem~4.4.2 and Chapter~1 Proposition~1.4.4,
Lemma~1.4.6, Lemma~1.4.12, Remark~1.4.14 and
Theorem~1.4.15 were freshly read, as were the norm/extension
Lemma~1.3.5 and the literal Galois-equivariance statement.
The Kummer/reflection argument needed here is given above,
with its sole root-of-unity correction explicitly removed
by ramification.  No whole-chapter read or independent proof
of deep class-field theory is claimed.
The integral Stickelberger theorem and its applied printed
proof retain the exact inverse group-ring convention of
(G168.8)--(G168.10).

The finite table has two independent exact implementations;
the residue translation is separately checked.
The original finite characters are unchanged, and the field
$H$ is used only to prove their residual strict vanishing.
All positive derived-index costs, the common window and its
simultaneous quantifiers, unlisted finite labels and coefficient
primes, the mixed input $\mathcal R_{\rm cub}$ and H04.4
remain separate.  No zero row in the finite table is presented
as a class-nonvanishing certificate.
